\clearpage
\phantomsection
\chapter*{Висновок}
\addcontentsline{toc}{chapter}{Висновок}
\label{chapter:conclusion}

У ході лабораторної роботи виконано попереднє проєктування гібридної RAG-системи для пошуку інформації по ігровій збірці GregTech: New Horizons (GTNH). Визначено цільових користувачів, типові завдання, критерії якості відповідей, словник термінів предметної області та джерела даних (офіційна вікі, книга квестів і база рецептів із nesql-exporter), обсяг яких обмежено приблизно до 10000 документів шляхом фільтрації й агрегації.

З урахуванням обмежень наявного обладнання (RTX 3060 12 ГБ, 32 ГБ ОЗП) обрано локальний стек: Python 3.13, PostgreSQL для точного пошуку рецептів, Qdrant із моделлю ембедингів BAAI/bge-m3 для текстового пошуку та Qwen 2.5 7B для генерації відповідей. Також побудовано загальну діаграму архітектури системи, що описує поєднання обох гілок пошуку в спільному контексті. Створено Github проєкту за посиланням \url{https://github.com/NureKostiukDenis/tops-labs}